\chapter{Teilchenklassifikation}

\section{Ziel dieses Kapitels}

Dieses Kapitel fuehrt die grundlegende Einteilung der Teilchenphysik ein.
Es geht darum, welche Teilchenarten es gibt, wie sie klassifiziert werden und welche
Wechselwirkungen fuer sie relevant sind.

\begin{notebox}
Dieses Kapitel behandelt die Stichwoerter:
\begin{itemize}
    \item Elementarteilchen,
    \item Fermionen,
    \item Bosonen,
    \item Leptonen,
    \item Quarks,
    \item Hadronen,
    \item Baryonen,
    \item Mesonen,
    \item Antiteilchen,
    \item Austauschteilchen,
    \item elektromagnetische Wechselwirkung,
    \item schwache Wechselwirkung,
    \item starke Wechselwirkung,
    \item Gravitation.
\end{itemize}
\end{notebox}

\section{Grundidee der Teilchenklassifikation}

In der Teilchenphysik will man Teilchen nicht nur aufzaehlen, sondern systematisch
ordnen.
Die Ordnung richtet sich nach Eigenschaften wie Masse, Spin, Ladung, Wechselwirkung
und innerer Struktur.

\begin{conceptbox}
Die wichtigsten Fragen bei der Klassifikation eines Teilchens sind:
\begin{itemize}
    \item Ist das Teilchen elementar oder zusammengesetzt?
    \item Hat es halbzahligen oder ganzzahligen Spin?
    \item Traegt es elektrische Ladung?
    \item Traegt es Farbladung?
    \item An welchen Wechselwirkungen nimmt es teil?
    \item Ist es ein Materieteilchen oder ein Austauschteilchen?
\end{itemize}
\end{conceptbox}

\begin{definitionbox}
Ein \emph{Elementarteilchen} ist ein Teilchen, das nach heutigem Wissen keine innere
Substruktur besitzt.
\end{definitionbox}

\begin{conceptbox}
Im Standardmodell sind Quarks, Leptonen und Eichbosonen elementar.

Hadronen wie Protonen, Neutronen, Pionen und Kaonen sind dagegen zusammengesetzte
Teilchen.
Sie bestehen aus Quarks und/oder Antiquarks.
\end{conceptbox}

\section{Fermionen und Bosonen}

Eine sehr grundlegende Einteilung erfolgt nach dem Spin.

\begin{definitionbox}
\emph{Fermionen} sind Teilchen mit halbzahligem Spin:
\[
    s=\frac{1}{2},\frac{3}{2},\ldots
\]

\emph{Bosonen} sind Teilchen mit ganzzahligem Spin:
\[
    s=0,1,2,\ldots
\]
\end{definitionbox}

\begin{conceptbox}
Materieteilchen sind im Standardmodell Fermionen.
Dazu gehoeren:
\begin{itemize}
    \item Quarks,
    \item Leptonen.
\end{itemize}

Austauschteilchen der Wechselwirkungen sind Bosonen.
Dazu gehoeren:
\begin{itemize}
    \item Photon,
    \item Gluonen,
    \item $W^+$, $W^-$ und $Z^0$,
    \item Higgs-Boson als skalares Boson.
\end{itemize}
\end{conceptbox}

\begin{examtipbox}
Merksatz:
Fermionen bilden Materie.
Bosonen vermitteln Wechselwirkungen.

Das ist eine sehr gute erste Orientierung, auch wenn es in Details Ausnahmen und
zusammengesetzte Bosonen gibt.
\end{examtipbox}

\section{Leptonen}

Leptonen sind elementare Fermionen, die keine starke Wechselwirkung besitzen.

\begin{definitionbox}
\emph{Leptonen} sind elementare Fermionen, die nicht an der starken Wechselwirkung
teilnehmen.
\end{definitionbox}

\begin{conceptbox}
Es gibt drei geladene Leptonen:
\[
    e^-,
    \qquad
    \mu^-,
    \qquad
    \tau^-.
\]

Dazu gehoeren drei Neutrinos:
\[
    \nu_e,
    \qquad
    \nu_\mu,
    \qquad
    \nu_\tau.
\]
\end{conceptbox}

\begin{notebox}
Die geladenen Leptonen besitzen elektrische Ladung
\[
    Q=-1.
\]
Die Neutrinos sind elektrisch neutral:
\[
    Q=0.
\]
\end{notebox}

\begin{conceptbox}
Leptonen nehmen an der schwachen Wechselwirkung teil.
Geladene Leptonen nehmen zusaetzlich an der elektromagnetischen Wechselwirkung teil.

Neutrinos nehmen nicht an der elektromagnetischen Wechselwirkung teil, weil sie
elektrisch neutral sind.
Sie werden daher hauptsaechlich ueber die schwache Wechselwirkung nachgewiesen.
\end{conceptbox}

\begin{examtipbox}
Elektron, Myon und Tauon sind geladene Leptonen.

Neutrinos sind neutrale Leptonen.

Leptonen sind keine Hadronen und bestehen nicht aus Quarks.
\end{examtipbox}

\section{Leptonfamilien}

Die Leptonen treten in drei Familien auf.

\begin{center}
\begin{tabular}{|c|c|c|}
\hline
Generation & geladenes Lepton & Neutrino \\
\hline
1 & $e^-$ & $\nu_e$ \\
\hline
2 & $\mu^-$ & $\nu_\mu$ \\
\hline
3 & $\tau^-$ & $\nu_\tau$ \\
\hline
\end{tabular}
\end{center}

\begin{conceptbox}
Die drei geladenen Leptonen haben dieselbe elektrische Ladung, unterscheiden sich
aber stark in ihrer Masse.

Das Elektron ist stabil.
Myon und Tauon sind instabil und zerfallen ueber die schwache Wechselwirkung.
\end{conceptbox}

\begin{notebox}
Zu jedem Lepton gibt es ein Antilepton:
\[
    e^+,
    \qquad
    \mu^+,
    \qquad
    \tau^+,
    \qquad
    \overline{\nu}_e,
    \qquad
    \overline{\nu}_\mu,
    \qquad
    \overline{\nu}_\tau.
\]
\end{notebox}

\section{Quarks}

Quarks sind elementare Fermionen mit Farbladung.
Sie bilden Hadronen.

\begin{definitionbox}
\emph{Quarks} sind elementare Fermionen, die Farbladung tragen und an der starken
Wechselwirkung teilnehmen.
\end{definitionbox}

\begin{conceptbox}
Es gibt sechs Quark-Flavors:
\[
    u,\quad d,\quad c,\quad s,\quad t,\quad b.
\]

Die Namen sind:
\begin{itemize}
    \item up,
    \item down,
    \item charm,
    \item strange,
    \item top,
    \item bottom.
\end{itemize}
\end{conceptbox}

\begin{notebox}
Quarks tragen elektrische Ladungen in Dritteln der Elementarladung.
Up-artige Quarks haben Ladung
\[
    Q=+\frac{2}{3}.
\]
Down-artige Quarks haben Ladung
\[
    Q=-\frac{1}{3}.
\]
\end{notebox}

\begin{center}
\begin{tabular}{|c|c|c|}
\hline
Generation & up-artiges Quark & down-artiges Quark \\
\hline
1 & $u$ mit $Q=+\frac{2}{3}$ & $d$ mit $Q=-\frac{1}{3}$ \\
\hline
2 & $c$ mit $Q=+\frac{2}{3}$ & $s$ mit $Q=-\frac{1}{3}$ \\
\hline
3 & $t$ mit $Q=+\frac{2}{3}$ & $b$ mit $Q=-\frac{1}{3}$ \\
\hline
\end{tabular}
\end{center}

\begin{examtipbox}
Die Quarkladungen muss man sicher koennen:
\[
    u,c,t:\quad +\frac{2}{3},
\]
\[
    d,s,b:\quad -\frac{1}{3}.
\]
\end{examtipbox}

\section{Antiteilchen}

Zu jedem Teilchen gibt es ein Antiteilchen.

\begin{definitionbox}
Ein \emph{Antiteilchen} besitzt dieselbe Masse und denselben Spin wie das zugehoerige
Teilchen, aber entgegengesetzte additive Quantenzahlen.
\end{definitionbox}

\begin{conceptbox}
Beispiele:
\begin{itemize}
    \item Zum Elektron $e^-$ gehoert das Positron $e^+$.
    \item Zum Proton $p$ gehoert das Antiproton $\overline{p}$.
    \item Zum Quark $u$ gehoert das Antiquark $\overline{u}$.
    \item Zum Neutrino $\nu_e$ gehoert das Antineutrino $\overline{\nu}_e$.
\end{itemize}
\end{conceptbox}

\begin{notebox}
Bei Antiteilchen kehren sich zum Beispiel elektrische Ladung, Baryonenzahl,
Leptonenzahl und Flavor-Quantenzahlen um.

Masse und Lebensdauer bleiben gleich.
\end{notebox}

\begin{examtipbox}
Nicht jedes neutrale Teilchen ist automatisch sein eigenes Antiteilchen.
Man muss die inneren Quantenzahlen betrachten.
\end{examtipbox}

\section{Hadronen}

Hadronen sind zusammengesetzte Teilchen aus Quarks.

\begin{definitionbox}
\emph{Hadronen} sind Teilchen, die aus Quarks und Antiquarks aufgebaut sind und an
der starken Wechselwirkung teilnehmen.
\end{definitionbox}

\begin{conceptbox}
Hadronen sind nicht elementar.
Sie besitzen eine innere Struktur.

Die wichtigsten Gruppen sind:
\begin{itemize}
    \item Baryonen,
    \item Mesonen.
\end{itemize}
\end{conceptbox}

\begin{definitionbox}
\emph{Baryonen} bestehen im einfachen Quarkmodell aus drei Quarks.

\emph{Mesonen} bestehen im einfachen Quarkmodell aus einem Quark und einem Antiquark.
\end{definitionbox}

\begin{formulabox}
Baryon:
\[
    qqq.
\]

Meson:
\[
    q\overline{q}.
\]
\end{formulabox}

\begin{examtipbox}
Hadron bedeutet:
Teilchen nimmt an der starken Wechselwirkung teil und besteht aus Quarks.

Leptonen sind keine Hadronen.
\end{examtipbox}

\section{Baryonen}

Baryonen bestehen aus drei Quarks.

\begin{definitionbox}
Ein \emph{Baryon} ist ein Hadron mit Baryonenzahl
\[
    B=1.
\]
Im einfachen Quarkmodell besteht ein Baryon aus drei Quarks.
\end{definitionbox}

\begin{conceptbox}
Die wichtigsten Baryonen fuer die Kernphysik sind:
\[
    p,
    \qquad
    n.
\]

Das Proton hat Quarkinhalt:
\[
    p=uud.
\]

Das Neutron hat Quarkinhalt:
\[
    n=udd.
\]
\end{conceptbox}

\begin{examplebox}
Berechne die Protonladung aus dem Quarkinhalt:
\[
    p=uud.
\]
Dann:
\[
    Q_p
    =
    Q_u+Q_u+Q_d
    =
    \frac{2}{3}
    +
    \frac{2}{3}
    -
    \frac{1}{3}
    =
    1.
\]

Das Proton hat also Ladung
\[
    +e.
\]
\end{examplebox}

\begin{examplebox}
Berechne die Neutronladung aus dem Quarkinhalt:
\[
    n=udd.
\]
Dann:
\[
    Q_n
    =
    Q_u+Q_d+Q_d
    =
    \frac{2}{3}
    -
    \frac{1}{3}
    -
    \frac{1}{3}
    =
    0.
\]

Das Neutron ist elektrisch neutral.
\end{examplebox}

\begin{notebox}
Weitere wichtige Baryonen sind zum Beispiel:
\[
    \Delta\text{-Teilchen},
    \qquad
    \Lambda,
    \qquad
    \Sigma,
    \qquad
    \Xi,
    \qquad
    \Omega^-.
\]
Diese werden im Kapitel zu Multipletts und Quantenzahlen genauer behandelt.
\end{notebox}

\section{Mesonen}

Mesonen bestehen aus einem Quark und einem Antiquark.

\begin{definitionbox}
Ein \emph{Meson} ist ein Hadron mit Baryonenzahl
\[
    B=0.
\]
Im einfachen Quarkmodell besteht ein Meson aus einem Quark und einem Antiquark.
\end{definitionbox}

\begin{conceptbox}
Wichtige Mesonen sind:
\[
    \pi^+,\quad \pi^0,\quad \pi^-,
\]
\[
    K^+,\quad K^0,\quad \overline{K}^0,\quad K^-.
\]
\end{conceptbox}

\begin{examplebox}
Das positiv geladene Pion hat den Quarkinhalt:
\[
    \pi^+ = u\overline{d}.
\]
Die Ladung ist:
\[
    Q_{\pi^+}
    =
    Q_u + Q_{\overline{d}}
    =
    \frac{2}{3}
    +
    \frac{1}{3}
    =
    1.
\]
\end{examplebox}

\begin{examplebox}
Das negativ geladene Pion hat den Quarkinhalt:
\[
    \pi^- = d\overline{u}.
\]
Die Ladung ist:
\[
    Q_{\pi^-}
    =
    Q_d + Q_{\overline{u}}
    =
    -\frac{1}{3}
    -
    \frac{2}{3}
    =
    -1.
\]
\end{examplebox}

\begin{examtipbox}
Bei Mesonen immer beachten:
Ein Antiquark besitzt die entgegengesetzte elektrische Ladung des zugehoerigen Quarks.
\end{examtipbox}

\section{Baryonenzahl}

\begin{definitionbox}
Die \emph{Baryonenzahl} $B$ ist eine additive Quantenzahl.

Ein Quark hat:
\[
    B=\frac{1}{3}.
\]

Ein Antiquark hat:
\[
    B=-\frac{1}{3}.
\]
\end{definitionbox}

\begin{conceptbox}
Fuer ein Baryon aus drei Quarks gilt:
\[
    B=\frac{1}{3}+\frac{1}{3}+\frac{1}{3}=1.
\]

Fuer ein Meson aus Quark und Antiquark gilt:
\[
    B=\frac{1}{3}-\frac{1}{3}=0.
\]
\end{conceptbox}

\begin{examtipbox}
Baryonen haben Baryonenzahl $B=1$.

Antibaryonen haben Baryonenzahl $B=-1$.

Mesonen haben Baryonenzahl $B=0$.
\end{examtipbox}

\section{Austauschteilchen}

Wechselwirkungen werden in der Quantenfeldtheorie durch Austauschteilchen beschrieben.

\begin{definitionbox}
Ein \emph{Austauschteilchen} ist ein Boson, das eine fundamentale Wechselwirkung
vermittelt.
\end{definitionbox}

\begin{conceptbox}
Wichtige Austauschteilchen sind:
\begin{itemize}
    \item Photon fuer die elektromagnetische Wechselwirkung,
    \item Gluonen fuer die starke Wechselwirkung,
    \item $W^+$, $W^-$ und $Z^0$ fuer die schwache Wechselwirkung.
\end{itemize}
\end{conceptbox}

\begin{notebox}
In Feynman-Diagrammen werden verschiedene Austauschteilchen oft verschieden gezeichnet:
\begin{itemize}
    \item Photonen als Wellenlinien,
    \item Gluonen als Spirallinien,
    \item $W$- und $Z$-Bosonen als gestrichelte oder spezielle Linien.
\end{itemize}
\end{notebox}

\begin{examtipbox}
Zuordnung merken:
\[
    \gamma
    \quad\Rightarrow\quad
    \text{elektromagnetische Wechselwirkung}.
\]
\[
    g
    \quad\Rightarrow\quad
    \text{starke Wechselwirkung}.
\]
\[
    W^\pm,\;Z^0
    \quad\Rightarrow\quad
    \text{schwache Wechselwirkung}.
\]
\end{examtipbox}

\section{Die fundamentalen Wechselwirkungen}

\begin{center}
\begin{tabular}{|p{3.3cm}|p{3.4cm}|p{4.2cm}|p{3.6cm}|}
\hline
Wechselwirkung & Austauschteilchen & wirkt auf & typische Rolle \\
\hline
stark & Gluonen & Farbladung, Quarks und Gluonen & bindet Quarks in Hadronen \\
\hline
elektromagnetisch & Photon & elektrische Ladung & Licht, Coulombkraft, Streuung \\
\hline
schwach & $W^\pm$, $Z^0$ & Quarks und Leptonen & Beta-Zerfall, Neutrinos \\
\hline
Gravitation & hypothetisches Graviton & Masse/Energie & makroskopisch wichtig \\
\hline
\end{tabular}
\end{center}

\begin{conceptbox}
In der Teilchenphysik sind vor allem starke, elektromagnetische und schwache
Wechselwirkung wichtig.

Gravitation ist fuer einzelne Elementarteilchen normalerweise extrem schwach und
wird im Standardmodell der Teilchenphysik nicht als quantisierte Wechselwirkung
beschrieben.
\end{conceptbox}

\begin{examtipbox}
Welche Wechselwirkung beteiligt ist, erkennt man oft an:
\begin{itemize}
    \item den beteiligten Teilchen,
    \item den erhaltenen oder verletzten Quantenzahlen,
    \item der typischen Lebensdauer,
    \item dem Austauschteilchen,
    \item ob Neutrinos beteiligt sind.
\end{itemize}
\end{examtipbox}

\section{Welche Teilchen nehmen an welcher Wechselwirkung teil?}

\begin{conceptbox}
Quarks:
\begin{itemize}
    \item starke Wechselwirkung, weil sie Farbladung tragen,
    \item elektromagnetische Wechselwirkung, weil sie elektrische Ladung tragen,
    \item schwache Wechselwirkung.
\end{itemize}

Geladene Leptonen:
\begin{itemize}
    \item elektromagnetische Wechselwirkung,
    \item schwache Wechselwirkung.
\end{itemize}

Neutrinos:
\begin{itemize}
    \item schwache Wechselwirkung.
\end{itemize}
\end{conceptbox}

\begin{notebox}
Hadronen nehmen an der starken Wechselwirkung teil, weil sie aus Quarks aufgebaut
sind.

Leptonen nehmen nicht an der starken Wechselwirkung teil.
\end{notebox}

\section{Teilchen und Stabilitaet}

Viele Teilchen sind instabil und zerfallen.
Ob ein Zerfall moeglich ist, haengt von Erhaltungssaetzen und Massen ab.

\begin{conceptbox}
Ein Teilchen kann nur zerfallen, wenn:
\begin{itemize}
    \item die Endteilchen insgesamt leichter sind,
    \item alle relevanten Erhaltungssaetze erfuellt sind,
    \item eine Wechselwirkung existiert, die den Zerfall erlaubt.
\end{itemize}
\end{conceptbox}

\begin{notebox}
Typische stabile oder sehr langlebige Teilchen im Alltag sind:
\begin{itemize}
    \item Photon,
    \item Elektron,
    \item Proton,
    \item Neutrinos,
    \item gebundene Atomkerne, wenn sie stabil sind.
\end{itemize}
Das freie Neutron ist dagegen instabil und zerfaellt ueber die schwache Wechselwirkung.
\end{notebox}

\begin{examtipbox}
Ein Zerfall ist nicht nur eine Frage der Masse.
Auch Erhaltungssaetze und Wechselwirkungsart entscheiden, ob der Zerfall erlaubt ist.
\end{examtipbox}

\section{Teilchenfamilien im Standardmodell}

\begin{center}
\begin{tabular}{|c|c|c|c|}
\hline
Typ & 1. Generation & 2. Generation & 3. Generation \\
\hline
up-artiges Quark & $u$ & $c$ & $t$ \\
\hline
down-artiges Quark & $d$ & $s$ & $b$ \\
\hline
geladenes Lepton & $e^-$ & $\mu^-$ & $\tau^-$ \\
\hline
Neutrino & $\nu_e$ & $\nu_\mu$ & $\nu_\tau$ \\
\hline
\end{tabular}
\end{center}

\begin{conceptbox}
Die erste Generation bildet die gewoehnliche Materie:
\begin{itemize}
    \item Up- und Down-Quarks bilden Protonen und Neutronen.
    \item Elektronen bilden die Atomhuelle.
    \item Elektron-Neutrinos treten bei schwachen Prozessen auf.
\end{itemize}
\end{conceptbox}

\begin{notebox}
Die zweite und dritte Generation enthalten schwerere Kopien der ersten Generation.
Diese Teilchen sind instabil und zerfallen typischerweise in leichtere Teilchen.
\end{notebox}

\section{Zusammengesetzte Teilchen und Elementarteilchen}

\begin{center}
\begin{tabular}{|p{4.2cm}|p{4.8cm}|p{4.8cm}|}
\hline
Kategorie & elementar? & Beispiele \\
\hline
Leptonen & ja & $e^-$, $\mu^-$, $\tau^-$, Neutrinos \\
\hline
Quarks & ja & $u$, $d$, $s$, $c$, $b$, $t$ \\
\hline
Eichbosonen & ja & $\gamma$, $g$, $W^\pm$, $Z^0$ \\
\hline
Baryonen & nein & $p$, $n$, $\Delta$, $\Omega^-$ \\
\hline
Mesonen & nein & $\pi$, $K$, Quarkonia \\
\hline
Atomkerne & nein & ${}^{4}\mathrm{He}$, ${}^{12}\mathrm{C}$ \\
\hline
\end{tabular}
\end{center}

\begin{examtipbox}
Proton und Neutron sind keine Elementarteilchen.
Sie sind Baryonen und bestehen aus Quarks.

Elektron und Quarks gelten im Standardmodell als elementar.
\end{examtipbox}

\section{Teilchenklassifikation in Reaktionen}

Bei Teilchenreaktionen ist die Klassifikation wichtig, weil sie hilft, die erlaubten
Wechselwirkungen und Erhaltungssaetze zu erkennen.

\begin{examplebox}
Betrachte den Beta-Minus-Zerfall:
\[
    n \rightarrow p + e^- + \overline{\nu}_e.
\]

Klassifikation:
\begin{itemize}
    \item $n$ ist ein Baryon.
    \item $p$ ist ein Baryon.
    \item $e^-$ ist ein Lepton.
    \item $\overline{\nu}_e$ ist ein Antilepton.
\end{itemize}

Der Prozess ist schwach, weil ein Quark-Flavor geaendert wird:
\[
    d \rightarrow u + W^-.
\]
\end{examplebox}

\begin{examplebox}
Betrachte die elektromagnetische Annihilation:
\[
    e^- + e^+ \rightarrow \gamma + \gamma.
\]

Klassifikation:
\begin{itemize}
    \item $e^-$ ist ein Lepton.
    \item $e^+$ ist ein Antilepton.
    \item $\gamma$ ist das Photon.
\end{itemize}

Der Prozess ist elektromagnetisch.
\end{examplebox}

\section{Mini-Zusammenfassung}

\begin{notebox}
Die wichtigsten Punkte dieses Kapitels sind:
\begin{itemize}
    \item Fermionen haben halbzahligen Spin.
    \item Bosonen haben ganzzahligen Spin.
    \item Leptonen sind elementare Fermionen ohne starke Wechselwirkung.
    \item Quarks sind elementare Fermionen mit Farbladung.
    \item Hadronen bestehen aus Quarks und nehmen an der starken Wechselwirkung teil.
    \item Baryonen bestehen im einfachen Quarkmodell aus drei Quarks.
    \item Mesonen bestehen im einfachen Quarkmodell aus einem Quark und einem Antiquark.
    \item Proton:
    \[
        p=uud.
    \]
    \item Neutron:
    \[
        n=udd.
    \]
    \item Quarkladungen:
    \[
        u,c,t:\quad +\frac{2}{3},
        \qquad
        d,s,b:\quad -\frac{1}{3}.
    \]
    \item Das Photon vermittelt die elektromagnetische Wechselwirkung.
    \item Gluonen vermitteln die starke Wechselwirkung.
    \item $W^\pm$ und $Z^0$ vermitteln die schwache Wechselwirkung.
    \item Die erste Generation bildet die normale Materie.
\end{itemize}
\end{notebox}

\section{Pruefungscheck}

\begin{examtipbox}
Nach diesem Kapitel solltest du folgende Fragen beantworten koennen:
\begin{enumerate}
    \item Was ist ein Elementarteilchen?
    \item Was ist der Unterschied zwischen Fermionen und Bosonen?
    \item Welche Teilchen sind Leptonen?
    \item Welche Leptonfamilien gibt es?
    \item Welche sechs Quark-Flavors gibt es?
    \item Welche elektrischen Ladungen haben die Quarks?
    \item Was ist ein Antiteilchen?
    \item Was ist ein Hadron?
    \item Was ist ein Baryon?
    \item Was ist ein Meson?
    \item Was ist der Quarkinhalt von Proton und Neutron?
    \item Wie berechnet man die Ladung eines Hadrons aus den Quarkladungen?
    \item Was ist die Baryonenzahl?
    \item Welche Austauschteilchen gehoeren zu welcher Wechselwirkung?
    \item Warum nehmen Leptonen nicht an der starken Wechselwirkung teil?
    \item Warum nehmen Quarks an der starken Wechselwirkung teil?
    \item Welche Teilchen bilden die erste Generation?
    \item Warum sind Proton und Neutron keine Elementarteilchen?
    \item Woran erkennt man, ob eine Reaktion elektromagnetisch, schwach oder stark ist?
\end{enumerate}
\end{examtipbox}